% Transcription of Cayley Collected Papers, Vol. 12, pp. 201--250
% Papers 825--828
\documentclass[11pt,a4paper]{article}
\usepackage[utf8]{inputenc}
\usepackage[T1]{fontenc}
\usepackage[english]{babel}
\usepackage{amsmath,amssymb,amsthm}
\usepackage{geometry}
\usepackage{fancyhdr}
\usepackage{titlesec}
\usepackage{enumitem}
\usepackage{array}
\usepackage{booktabs}
\usepackage{longtable}
\usepackage{slashed}
\usepackage{mathrsfs}

\geometry{margin=1in}

\pagestyle{fancy}
\fancyhf{}
\fancyhead[C]{\small Cayley, Collected Math. Papers, Vol. XII, pp. 201--250}
\fancyfoot[C]{\thepage}
\renewcommand{\headrulewidth}{0.4pt}

\theoremstyle{definition}
\newtheorem{art}{Art.}[section]

\DeclareMathOperator{\sech}{sech}
\DeclareMathOperator{\cosec}{cosec}

\title{\textbf{Transcription}\\[6pt]
\Large Arthur Cayley: Collected Mathematical Papers\\
Volume XII, Pages 201--250\\[12pt]
\large Papers 825--828}
\author{Transcribed from the original collected works}
\date{}

\begin{document}
\maketitle
\thispagestyle{empty}

\tableofcontents
\newpage

%%%%%%%%%%%%%%%%%%%%%%%%%%%%%%%%%%%%%%%%%%%%%%%%%%%%%%%%%%%%%%%%%%%%%%%%%%%%%%%
\section{Paper 825. A Memoir on the Abelian and Theta Functions (pp. 201--215)}
%%%%%%%%%%%%%%%%%%%%%%%%%%%%%%%%%%%%%%%%%%%%%%%%%%%%%%%%%%%%%%%%%%%%%%%%%%%%%%%

\noindent
[From the \textit{Journal f\"ur die reine und angewandte Mathematik} (Crelle), t.~xciv. (1883), pp.~186--210.]

\subsection*{Application of Abel's Theorem. Art.~Nos. 153 to 157}

\begin{art}
Taking the fixed curve to be $f = (Az^2 + 2Bz + C) = 0$, we have
\[
\frac{dz}{y} = \frac{dy}{dx},
\]
if for shortness we write
\[
(x,y)^2 = B^2 - AC = (x-ay)(x-by)(x-cy)(x-dy)(x-ey)(x-fy),
\]
and we thence have
\[
\omega = \int \frac{x\,dy - y\,dx}{\sqrt{(x,y)^6}}.
\]

The minor curve $(x,y,z)^{n-3} = 0$ is an arbitrary line passing through the node, that is, the point $x=0$, $y=0$; and the pure theorem thus gives the two relations
\[
\Sigma x\,d\omega = 0, \qquad \Sigma y\,d\omega = 0;
\]
where the summation extends to the intersections of the fixed curve $Az^2 + 2Bz + C = 0$ with the variable curve $\phi$.

The variable curve is taken to be a cubic
\[
Az + B = (\alpha,\beta,\gamma,\delta\,\raisebox{0.5ex}{$\scriptstyle\diagup$}\,e,f)(x,y)^3,
\]
or say $Az + B = \Omega$, where $\Omega$ is a given cubic function $(x,y)^3$: viz. this is a nodal cubic, the node and the two tangents there being the same with the node and the two tangents of the quartic: hence it meets the quartic in the node counting 6 times, and in 6 other points, say these are the points $1,2,3,4,5,6$: hence the differential relations are
\begin{align*}
x_1\,d\omega_1 + x_2\,d\omega_2 + x_3\,d\omega_3 + x_4\,d\omega_4 + x_5\,d\omega_5 + x_6\,d\omega_6 &= 0, \\
y_1\,d\omega_1 + y_2\,d\omega_2 + y_3\,d\omega_3 + y_4\,d\omega_4 + y_5\,d\omega_5 + y_6\,d\omega_6 &= 0.
\end{align*}
\end{art}

\begin{art}
Observe that the intersections of the cubic with the fixed curve are given by the equation $H^2 = B^2 - AC$, or say $H^2 = (x,y)^6$, an equation which determines the ratio $x:y$ for the six points respectively; and the ratio $z:x$ is then determined rationally by the original equation $Az + B = \Omega$.

Instead of regarding $\Omega$ as a given function, we may, if we please, take $1,2,3,4$ given points on the quartic: we then have four equations for the determination of the coefficients $(\alpha,\beta,\gamma,\delta)$ of the function $\Omega$; viz. these equations may be taken to be
\[
(\alpha,\beta,\gamma,\delta)(x_1,y_1)^3 = \Omega_1, \quad (\alpha,\beta,\gamma,\delta)(x_2,y_2)^3 = \Omega_2, \quad \text{etc.},
\]
$\Omega$ is hereby completely determined: and this being so, the remaining intersections $5$ and $6$ are also completely determined: there are thus between the six points 2 integral relations, which agrees with the number, $=2$, of the differential relations obtained above.
\end{art}

\begin{art}\label{art:155}
If we now assume
\begin{align*}
du &= x_1\,d\omega_1 + x_2\,d\omega_2, &\qquad dv &= y_1\,d\omega_1 + y_2\,d\omega_2, \\
du' &= x_3\,d\omega_3 + x_4\,d\omega_4, &\qquad dv' &= y_3\,d\omega_3 + y_4\,d\omega_4, \\
du'' &= x_5\,d\omega_5 + x_6\,d\omega_6, &\qquad dv'' &= y_5\,d\omega_5 + y_6\,d\omega_6,
\end{align*}
or say
\begin{align*}
u,v &= \frac{\int x(x\,dy - y\,dx)}{\sqrt{(x,y)^6}}, &\qquad
u',v' &= \frac{\int x(x\,dy - y\,dx)}{\sqrt{(x,y)^6}}, \\
u'',v'' &= \frac{\int x(x\,dy - y\,dx)}{\sqrt{(x,y)^6}}, &\qquad
v &= \frac{\int y(x\,dy - y\,dx)}{\sqrt{(x,y)^6}},
\end{align*}
that is,
\[
u = \int_{\alpha}^{\beta} \frac{x(x\,dy - y\,dx)}{\sqrt{(x,y)^6}}, \qquad
v = \int_{\alpha}^{\beta} \frac{y(x\,dy - y\,dx)}{\sqrt{(x,y)^6}},
\]
where $\alpha$, $\beta$ are points assumed at pleasure on the quartic: and similarly for $u'$, $v'$: $u''$, $v''$: then $u$, $v$ are hereby determined as functions of the points $1,2$: and we may conversely regard the points $1,2$ as determined in terms of the two arguments $u$, $v$.

We might, selecting any two symmetrical functions of the degree zero, for instance, $\frac{x_1}{y_1} + \frac{x_2}{y_2}$, $\frac{x_1}{y_1} \cdot \frac{x_2}{y_2}$, represent them as functions $\phi(u,v)$, $\psi(u,v)$: and then $\frac{x_1}{y_1}$ and $\frac{x_2}{y_2}$ will be functions of $\phi(u,v)$, $\psi(u,v)$. But instead of this selection, it is proper to consider the ratios of six functions depending on the points $a,b,c,d,e,f$ respectively: viz. we assume
\begin{multline*}
\sqrt{(x_1-ay_1)(x_2-ay_2)} : \sqrt{(x_1-by_1)(x_2-by_2)} : \cdots : \sqrt{(x_1-fy_1)(x_2-fy_2)} \\
= A(u,v) : B(u,v) : \cdots : F(u,v),
\end{multline*}
and of course $3,4$ will be in like manner determined by means of the corresponding functions of $u'$, $v'$, and $5,6$ by means of the corresponding functions of $u''$, $v''$.

The squared functions $A^2$, $B^2$, $C^2$, $D^2$, $E^2$, $F^2$ are proportional to given linear functions of $x_1x_2$, $x_1y_2+x_2y_1$, $y_1y_2$, and are thus connected by three independent linear relations.
\end{art}

\begin{art}
The differential relations then become
\[
du + du' + du'' = 0, \qquad dv + dv' + dv'' = 0,
\]
and we have consequently
\[
u + u' + u'' = I, \qquad v + v' + v'' = J,
\]
where $I$, $J$ are constants which are determinable as definite integrals by the consideration that, when the cubic is taken to be $Az + B = 0$, the six points $1,2,3,4,5,6$ coincide with the points of contact $a,b,c,d,e,f$.

I do not at present see my way to a proper development of this point of the theory: but in explanation of the nature of the result, I assume for the moment that by a proper determination of the inferior limits $\alpha$, $\beta$, or otherwise, we may take $I=0$, $J=0$.

We then have $u'' = -u - u'$, $v'' = -v - v'$: and the integral equations, which determine the points $5,6$ in terms of the points $1,2$ and the points $3,4$, then in effect determine the functions $A$, $B$, \&c., of $-u-u'$, $-v-v'$, or say those of $u+u'$, $v+v'$ in terms of the like functions of $(u,v)$ and of $(u',v')$: viz. these equations give the addition-theory of the functions $A(u,v)$, \&c.
\end{art}

\begin{art}
We may, in the first instance, disregarding altogether the consideration of the arguments $u$, $v$, \&c., attend only to the algebraic functions such as $\sqrt{(x_1-ay_1)(x_2-ay_2)}$, \&c., of the coordinates of the pairs of points $1,2$; $3,4$, and $5,6$; and we can in regard to these develope a proper theory. This depends only on the equation $H = \sqrt{(x,y)^6}$; it will be convenient to assume herein $y=1$, and slightly modifying the form, to write it
\[
(\alpha,\beta,\gamma,\delta)(x,1)^3 = \sqrt{a-x \cdot b-x \cdot c-x \cdot d-x \cdot e-x \cdot f-x};
\]
and accordingly to consider the functions $\sqrt{a-x_1 \cdot a-x_2}$, \&c. These are called the single-letter functions $A$, \&c., but there are certain double-letter functions $AB$, \&c., which have also to be considered; and I will, in the first instance, show how these present themselves in connexion with the cubic curve.
\end{art}

\subsection*{Origin of the Double-Letter Functions. Art.~Nos. 158 and 159}

\begin{art}
The cubic curve $Az + B = \Omega$ may be taken to be a curve through two of the points of contact, say the points $a,b$; these will then be two out of the six points, and taking the remaining four points to be the pairs $1,2$ and $3,4$, we have single-letter functions of $3,4$ presenting themselves as double-letter functions of $1,2$.

In fact, the equation of the curve is
\[
Az + B = \lambda (x-ay)(x-by)(x-ky);
\]
for the intersections with the quartic we have
\[
\lambda^2 (x-ay)^2 (x-by)^2 (x-ky)^2 = \Omega^2,
\]
or throwing out the factor $(x-ay)(x-by)$ and changing the constant $\lambda$, this is
\[
(x-ay)(x-by)(x-ky)^2 - \mu (x-cy)(x-dy)(x-ey)(x-fy) = 0;
\]
and the quartic function must be a multiple of $(xy_1-x_1y)(xy_2-x_2y)(xy_3-x_3y)(xy_4-x_4y)$.

Putting each of the $y$'s equal $1$, we have the identity
\[
(a-x)(b-x)(k-x)^2 - \mu(c-x)(d-x)(e-x)(f-x) = \rho(x_1-x)(x_2-x)(x_3-x)(x_4-x);
\]
and hence, introducing a notation which will be convenient, $a-x = a_x$, $a-x_1 = a_1$, and so in other cases, we have by giving different values to $x$ the equations
\[
\begin{aligned}
a_2b_2k_2^{\,2} &= \mu c_2d_2e_2f_2, &\qquad (a-c)(b-c)(k-c)^2 &= \mu c_1c_2c_3c_4, \\
a_3b_3k_3^{\,2} &= \mu c_3d_3e_3f_3, &\qquad (a-d)(b-d)(k-d)^2 &= \mu d_1d_2d_3d_4, \\
a_4b_4k_4^{\,2} &= \mu c_4d_4e_4f_4, &\qquad (a-e)(b-e)(k-e)^2 &= \mu e_1e_2e_3e_4, \\[-2pt]
& & (a-f)(b-f)(k-f)^2 &= \mu f_1f_2f_3f_4,
\end{aligned}
\]
and
\[
\lambda(c-a)(d-a)(e-a)(f-a) = \mu a_1a_2a_3a_4, \qquad
\lambda(c-b)(d-b)(e-b)(f-b) = \mu b_1b_2b_3b_4.
\]
\end{art}

\newpage
\begin{art}
We have thus
\[
\frac{(a-c)(b-c)}{c_1c_2 \cdot c_3c_4} = \frac{(a-d)(b-d)}{d_1d_2 \cdot d_3d_4},
\]
and
\[
\frac{k_1k_2}{a_2b_2 \cdot c_1d_1e_1f_1} = \frac{k_3k_4}{a_3b_3 \cdot c_3d_3e_3f_3},
\]
which last equation, writing for a moment
\[
\gamma = \sqrt{a_2b_2 \cdot c_1d_1e_1f_1}, \qquad
\delta = \sqrt{a_1b_1 \cdot c_2d_2e_2f_2},
\]
gives
\[
k(\gamma-\delta) = \#\gamma - \#\delta, \quad \text{and thence} \quad
\frac{\gamma}{\delta} = \frac{k-\delta}{k-\gamma},
\]
\[
\frac{\gamma c_3c_4}{\delta d_3d_4} = \frac{\sqrt{a_2b_2c_1d_1e_1f_1} \cdot c_3c_4}{\sqrt{a_1b_1c_2d_2e_2f_2} \cdot d_3d_4},
\]
or, substituting in the first equation, we have
\[
\frac{\sqrt{(a-c)(b-c)}}{\sqrt{(a-d)(b-d)}} = \frac{\sqrt{a_2b_2c_1d_1e_1f_1} \cdot c_3c_4}{\sqrt{a_1b_1c_2d_2e_2f_2} \cdot d_3d_4}.
\]
\end{art}

\begin{art}
Considering the duad $DE$ as an abbreviation for the double triad $ABC.DEF$, the expressed duad being always accompanied by the letter $F$, we are thus led to the consideration of the double-letter functions
\[
ABK = \frac{1}{x_1-x_2}\left\{\sqrt{a_1b_1f_1 \cdot c_2d_2e_2} - \sqrt{a_2b_2f_2 \cdot c_1d_1e_1}\right\}, \quad \text{\&c.},
\]
in connexion with the already mentioned single-letter functions
\[
A_{12} = \sqrt{a_1a_2}, \quad \text{\&c.},
\]
viz. in this notation the equation just obtained is
\[
\frac{C_3C_4}{D_3D_4} = \frac{\sqrt{(a-c)(b-c)}}{\sqrt{(a-d)(b-d)}} \cdot \frac{DEK}{CEF},
\]
and it thus appears that, the points $3,4$ being obtained as above from the given points $1,2$, then the quotient of two of the single-letter functions of $3,4$ is a constant multiple of the quotient of two of the double-letter functions of $1,2$.

Observe that the points $3,4$ are derived from $1,2$ by means of the two points $a,b$: we have $DE$ standing for $ABC.DEF$, $CE$ for $ABD.CEF$, and if the two functions were represented by $ABC$, $ABD$ respectively, then the form would have been
\[
\frac{\sqrt{(a-c)(b-c)}}{\sqrt{(a-d)(b-d)}} \cdot \frac{ABC}{ABD},
\]
which is a clearer expression of the theorem; the apparent want of symmetry of the first form arises only from the arbitrary selection of the letter $F$ to accompany the expressed duad, and is at once removed by substituting for a duad $DE$ the triad $ABC.DEF$ which is thereby signified.

The denominator factor $x_1-x_2$ is introduced in order to make the degree in $x_1$ or $x_2$ equal to that of the single-letter functions.
\end{art}

\newpage
\subsection*{The Addition Theory. Art.~Nos. 160 to 163}

\begin{art}
We have the six single-letter symbols $A$, $B$, $C$, $D$, $E$, $F$; viz. $A_{12}^2 = a_1a_2$, \&c.: and the ten double-letter symbols $AB$, $AC$, $AD$, $AE$, $BC$, $BD$, $BE$, $CD$, $CE$, $DE$, viz.
\[
ABK = \frac{1}{x_1-x_2}\left\{\sqrt{a_1b_1f_1 \cdot c_2d_2e_2} - \sqrt{a_2b_2f_2 \cdot c_1d_1e_1}\right\}, \quad \text{\&c.},
\]
these 16 functions being connected by algebraical relations which are immediately deducible from these expressions of the functions in terms of $x_1$, $x_2$.

The problem is to express the functions of $5,6$ in terms of those of $1,2$ and of those of $3,4$.

The relation between the variables $x_1$, $x_2$, $x_3$, $x_4$, $x_5$, $x_6$ consists herein that we have $x_1$, $x_2$, $x_3$, $x_4$, $x_5$, $x_6$ as the roots of the equation
\[
(\alpha x^3 + \beta x^2 + \gamma x + \delta)^2 - \lambda(a-x)(b-x)(c-x)(d-x)(e-x)(f-x) = 0;
\]
or, what is the same thing, it consists in the identity
\begin{multline*}
(\alpha x^3 + \beta x^2 + \gamma x + \delta)^2 - \lambda(a-x)(b-x)(c-x)(d-x)(e-x)(f-x) \\
= \rho(x-x_1)(x-x_2)(x-x_3)(x-x_4)(x-x_5)(x-x_6).
\end{multline*}
Again, it may be expressed by the plexus of equations
\[
\frac{1}{x_1},\; \frac{1}{x_2},\; \frac{1}{x_3},\; \frac{1}{x_4},\; \frac{1}{x_5},\; \frac{1}{x_6}, \quad
\frac{1}{x_1^2},\; \frac{1}{x_2^2},\; \frac{1}{x_3^2},\; \frac{1}{x_4^2},\; \frac{1}{x_5^2},\; \frac{1}{x_6^2}, \quad \&\text{c.},
\]
where $X_1 = (a-x_1)(b-x_1)\cdots(f-x_1)$, \&c.; these are equivalent of course to two equations, and serve to determine $x_5$, $x_6$ in terms of $x_1$, $x_2$, $x_3$, $x_4$.
\end{art}

\begin{art}
The solution is, in fact, as is given in my paper ``On the addition of the double $\mathcal{P}$-functions,'' Crelle, t.~LXXXVIII. (1880), pp.~74--81, [703].

Writing successively $x = x_1$, $x_2$, $x_3$, $x_4$, we have
\[
\alpha x_1^3 + \beta x_1^2 + \gamma x_1 + \delta = \sqrt{\mu \cdot A_1 B_1 C_1 D_1 E_1 F_1},
\]
which equations serve to determine the ratios $\alpha:\beta:\gamma:\delta$ in terms of $x_1$, $x_2$, $x_3$, $x_4$; and we have then the two like equations
\[
\alpha x_5^3 + \beta x_5^2 + \gamma x_5 + \delta = \sqrt{\mu \cdot A_5 B_5 C_5 D_5 E_5 F_5}, \qquad
\alpha x_6^3 + \beta x_6^2 + \gamma x_6 + \delta = \sqrt{\mu \cdot A_6 B_6 C_6 D_6 E_6 F_6},
\]
which determine the symmetric functions of $x_5$, $x_6$.
\end{art}

\begin{art}
If, reverting to the identity, we write therein for instance $x = a$, we find
\[
\alpha a^3 + \beta a^2 + \gamma a + \delta = \sqrt{\mu \cdot A_1 A_2 A_3 A_4 A_5 A_6},
\]
which equation when properly reduced gives the proportional value of $A_{56}$.
\end{art}

\begin{art}
Calling for a moment the function on the left-hand side $\Omega$, we have
\[
\begin{vmatrix}
x_1^3 & x_1^2 & x_1 & 1 & \sqrt{X_1} \\
x_2^3 & x_2^2 & x_2 & 1 & \sqrt{X_2} \\
x_3^3 & x_3^2 & x_3 & 1 & \sqrt{X_3} \\
x_4^3 & x_4^2 & x_4 & 1 & \sqrt{X_4} \\
a^3 & a^2 & a & 1 & \sqrt{A_1 A_2 A_3 A_4 A_5 A_6}
\end{vmatrix} = 0,
\]
that is,
\[
\begin{vmatrix}
x_1^3 & x_1^2 & x_1 & 1 & \sqrt{X_1} \\
x_2^3 & x_2^2 & x_2 & 1 & \sqrt{X_2} \\
x_3^3 & x_3^2 & x_3 & 1 & \sqrt{X_3} \\
x_4^3 & x_4^2 & x_4 & 1 & \sqrt{X_4} \\
a^3 & a^2 & a & 1 & A_{56}
\end{vmatrix} = 0,
\]
viz. this is
\[
\begin{vmatrix}
x_1^3 & x_1^2 & x_1 & 1 & \sqrt{X_1} \\
x_2^3 & x_2^2 & x_2 & 1 & \sqrt{X_2} \\
x_3^3 & x_3^2 & x_3 & 1 & \sqrt{X_3} \\
x_4^3 & x_4^2 & x_4 & 1 & \sqrt{X_4} \\
a^3 & a^2 & a & 1 & A_{56}
\end{vmatrix} = 0,
\]
or, as this may be written,
\[
\sqrt{X_3 \cdot a - x_3 \cdot a - x_4} \cdot \frac{x_4 - a}{x_2 - x_1 \cdot x_4 - x_3} + \cdots = 0.
\]

We have here $\sqrt{a - x_3 \cdot a - x_4} = \sqrt{A_{34}}$, and the function which multiplies this is without difficulty found to be
\[
\nabla_A = \partial_b + \partial_c + \partial_d,
\]
where the summation extends to the three terms obtained by the cyclical interchange of the letters $b,c,d$: these being a set of three out of the five letters other than $a$.

Similarly $\sqrt{a - x_1 \cdot a - x_2} = \sqrt{A_{12}}$, and the function which multiplies this is
\[
\frac{c-d}{b-c} \cdot \frac{d}{b} \cdot \frac{b}{c}.
\]

The expression for $\Omega$ thus contains the factor $A_{12} A_{34}$: but we have this equation contains therefore the factor $\sqrt{A_{12} A_{34}}$, and omitting it we find
\[
\frac{1}{\sqrt{?}} (x_1 - x_3 \cdot x_1 - x_4 \cdot x_2 - x_3 \cdot x_2 - x_4)(c-d \cdot d-b \cdot b-c) A_{56} + \cdots = 0,
\]
where, as before, the summations refer each to the three terms obtained by the cyclical interchange of the letters $b,c,d$; these being any three of the five letters other than $a$: and the remaining two letters $e,f$ enter into the formulae symmetrically.

The formula thus gives for $A_{56}$ ten values which are of course equal to each other. By reason of the undetermined factor $\frac{1}{\sqrt{?}}$, the formula gives only the proportional value of $A_{56}$; viz. combining it with the like formulae for $B_{56}$, \&c., we have determinate values of the ratios $A_{56} : B_{56} : \cdots : F_{56}$. But this being understood, we regard the formula as a formula for each single-letter function of $x_5$, $x_6$ in terms of the single and double-letter functions of $x_1$, $x_2$ and of $x_3$, $x_4$ respectively.
\end{art}

\begin{art}
We require further the expressions for the double-letter functions of $5,6$. Consider for example the function $DE_{56}$, which is
\[
\frac{1}{x_5 - x_6}\left\{\sqrt{d_5 e_5 f_5 \cdot a_6 b_6 c_6} - \sqrt{d_6 e_6 f_6 \cdot a_5 b_5 c_5}\right\};
\]
then multiplying by $A_{56} B_{56} C_{56} = \sqrt{a_5 b_5 c_5 \cdot a_6 b_6 c_6}$, we have
\[
\frac{1}{x_5 - x_6}\{a_6 b_6 c_6 \sqrt{X_5} - a_5 b_5 c_5 \sqrt{X_6}\},
\]
or recollecting that $\sqrt{X_5}$ and $\sqrt{X_6}$ are $= \alpha x_5^3 + \beta x_5^2 + \gamma x_5 + \delta$ and $\alpha x_6^3 + \beta x_6^2 + \gamma x_6 + \delta$ respectively, this may be written
\begin{multline*}
\frac{1}{x_5 - x_6}\bigl\{(a-x_5)(b-x_5)(c-x_5)(\alpha x_6^3 + \beta x_6^2 + \gamma x_6 + \delta) \\
- (a-x_6)(b-x_6)(c-x_6)(\alpha x_5^3 + \beta x_5^2 + \gamma x_5 + \delta)\bigr\}.
\end{multline*}

Using the well-known identity
\begin{multline*}
\frac{1}{x_5 - x_6}\bigl\{(a-x_5)(b-x_5)(c-x_5)(\alpha x_6^3 + \cdots) \\
- (a-x_6)(b-x_6)(c-x_6)(\alpha x_5^3 + \cdots)\bigr\} = \sum \frac{(b-a)(c-a)(d-a)}{\sqrt{\cdots}},
\end{multline*}
where the summation extends to the four terms obtained by the cyclical interchanges of the letters $a,b,c,d$: and the like identity for $\alpha x_5^3 + \beta x_5^2 + \gamma x_5 + \delta$: there will be terms in $\alpha a^3 + \beta a^2 + \gamma a + \delta$, $\alpha b^3 + \beta b^2 + \gamma b + \delta$, $\alpha c^3 + \beta c^2 + \gamma c + \delta$, but the term in $\alpha d^3 + \beta d^2 + \gamma d + \delta$ will disappear of itself.

After some easy reductions, the formula for $DE_{56}$ is obtained in the form
\[
DE_{56} = \frac{\sqrt{A_{12} B_{34} C_{34}}}{\sqrt{D_{12} E_{12}}} \cdot \frac{(a-d)(b-d)}{(c-d)} + \text{symmetric terms}.
\]
\end{art}

\newpage
\subsection*{Chapter VII. The Functions T, U, V, $\omega$}

The present chapter is substantially a reproduction of C. and G.'s seventh section, ``Die Function $T_{12}$'', borrowing only from the next section the definition of the theta-function; but for greater completeness I give some developments which are not contained in the original memoir.

\begin{art}
Suppose that the curves $\phi$, $\psi$ are each of them a major curve, that is, a curve of the order $n-2$ passing through the $\delta$ dps, and consequently besides meeting the curve $f$ in $n(n-2)-2\delta = 2p+n-2$ points: the theorem is that if $\phi$, $\psi$ each pass through the same $(2p+n-2)-m$, $=q-m$ points respectively, then the remaining $m$ points of intersection with $f$ of the curve $\phi$ and the remaining $m$ points of intersection with $f$ of the curve $\psi$ are pairs of points belonging to the same $g_{m}^{m-1}$; that is, we have a relation $T(1,2,\ldots,m) = 0$, where $T$ is a symmetrical function of the $m$ points $1,2,\ldots,m$, and of the parametric points and of the sets of superior and inferior points respectively.
\end{art}

\begin{art}
The foregoing theorem for the quartic thus is
\[
T(a,b,c; a_1,b_1,c_1) = 0.
\]
Observing that
\[
\frac{T_a + T_b + T_c}{3} = 2\Sigma - \frac{T_{a_1} + T_{b_1} + T_{c_1}}{3},
\]
where $\Sigma$ denotes the sum of the parametric terms, we have
\[
T(a,b,c; a_1,b_1,c_1) = \int_a^{a_1} + \int_b^{b_1} + \int_c^{c_1}.
\]
\end{art}

\begin{art}
We have that is,
\[
T_{123} = \Sigma\omega - \Sigma\omega(4,5,6),
\]
and thence the theorem (A):
\[
T_{123} + T_{145} + T_{167} + T_{189} = \text{const}.
\]
\end{art}

\begin{art}
Theorem (B). We have
\[
T_{12}(\xi,\eta,\zeta) + T_{12}(\xi',\eta',\zeta') = 0,
\]
where $\xi,\eta,\zeta$ and $\xi',\eta',\zeta'$ are comajors in regard to $1,2$.
\end{art}

\begin{art}
To show that $T - T_1$ contains no term in $\log \mu$. For $T$, the only term in $\log \mu$ is $\int \frac{d\omega}{\mu}$, and for $T_1$ it is $\int \frac{d\omega_1}{\mu_1}$; and it is to be shown that the difference of the two integrals contains no term in $\log \mu$. Considering on the quartic the two new points $8,9$, we have $T_{12}(8,9,\mu) = T_{12}(8,9,\mu_1)$; and the theorem is thus proved.
\end{art}

\begin{art}
Theorem (C). We have
\[
T_{123}(\xi,\eta,\zeta) + T_{123}(\xi',\eta',\zeta') = 0,
\]
tag{C}
where $\xi,\eta,\zeta$ and $\xi',\eta',\zeta'$ are comajors in regard to $1,2,3$: the equation shows that, in $T(1,2,3)$ the function $T_{123}$, we can without alteration of the value interchange a pair of points $1$, $1'$ out of the system of comajor points.
\end{art}

\begin{art}
Taking the theorem for the quartic in the form $T_{123} + T_{145} + \text{const} = 0$, and applying it successively to the parametric line, and to the conic which determines the points $1,2,3$, we have
\[
P_1 + P_2 + P_3 + 2\Sigma = \text{const}, \qquad P_1' + P_2' + P_3' + 2\Sigma = \text{const}.
\]
Taking the difference of these equations, we have
\[
P_1 + P_2 + P_3 - P_1' - P_2' - P_3' = \text{const},
\]
viz. the points $\eta$, $1'$, $2'$, $3'$ being fixed points, this is
\[
P_1 + P_2 + P_3 = \text{const.},
\]
that is, the functions $u$, $v$, $w$ defined as above are each of them constant under the variation in question.
\end{art}

\newpage
%%%%%%%%%%%%%%%%%%%%%%%%%%%%%%%%%%%%%%%%%%%%%%%%%%%%%%%%%%%%%%%%%%%%%%%%%%%%%%%
\section{Paper 826. Note on a Partition-Series (pp. 216--219)}
%%%%%%%%%%%%%%%%%%%%%%%%%%%%%%%%%%%%%%%%%%%%%%%%%%%%%%%%%%%%%%%%%%%%%%%%%%%%%%%

\noindent
[From the \textit{American Journal of Mathematics}, vol.~vi. (1884), pp.~63, 64.]

Prof.~Sylvester, in his paper, ``A Constructive theory of Partitions, \&c.,'' \textit{Amer.\ Journ.\ Math.}, vol.~v. (1883), p.~282, has given the following very beautiful formula
\begin{multline*}
(1+ax)(1+ax^2)(1+ax^3)\cdots \\
= 1 + \frac{1+ax}{1}\,xa + \frac{(1+ax)(1+ax^2)}{1\cdot 2}\,x^3a^2 + \frac{(1+ax)(1+ax^2)(1+ax^3)}{1\cdot 2\cdot 3}\,x^6a^3 + \cdots,
\end{multline*}
or, as this may be written,
\[
\Omega = 1 + P + Q(1+ax) + R(1+ax)(1+ax^2) + S(1+ax)(1+ax^2)(1+ax^3) + \cdots,
\]
where the law is obvious.

I propose to show how the several coefficients $P$, $Q$, $R$, \&c. may be obtained in succession.

Writing for shortness $1,2,3,\ldots$ to denote $xa$, $x^3a^2$, $x^6a^3$, \&c. respectively, and $P$, $Q$, $R$, \&c. to denote the coefficients, we have
\[
\Omega = 1 + P + Q(1+ax) + R(1+ax)(1+ax^2) + S(1+ax)(1+ax^2)(1+ax^3) + \cdots,
\]
and we find without difficulty (see infra) that
\[
1 + Q + R = (1+ax)(1+R), \qquad 1 + R + S = (1+ax^2)(1+S), \quad \text{\&c.};
\]
and hence, using $\Pi$ to denote the sum
\[
\Omega = 1 + P + Q(1+ax) + R(1+ax)(1+ax^2) + S(1+ax)(1+ax^2)(1+ax^3) + \cdots,
\]
we obtain successively
\begin{align*}
\Omega - (1+ax) &= 1 + P + Q + R(1+ax^2) + S(1+ax^2)(1+ax^3) + \cdots, \\
\Omega - (1+ax)(1+ax^2) &= 1 + Q + R + S(1+ax^3) + \cdots,
\end{align*}
and so on.

In these equations, on the right-hand sides, the terms are successively from the beginning $1+P$, $1+Q$, $1+R$, $1+S$, \&c.; and we thus have
\[
\Omega = 1 + P + Q(1+ax) + R(1+ax)(1+ax^2) + S(1+ax)(1+ax^2)(1+ax^3) + \cdots,
\]
and therefore
\[
1 \cdot 2 \cdot 3 \cdot 4 \cdots = 1 - (1+x) + (1+x^2) - (1+x^3) + \cdots,
\]
which is Euler's theorem.

It might appear that the identities used in the proof would also, for this particular value $a=1$, lead to interesting theorems; but this is found not to be the case: we have
\[
P = \frac{x}{1}, \qquad Q = \frac{x^3}{1\cdot 2}, \qquad R = \frac{x^6}{1\cdot 2\cdot 3},
\]
but the expressions in terms of these quantities for the products $2\cdot 3\cdot 4\cdots$, $3\cdot 4\cdots$, \&c., contain denominator factors, and are thus altogether without interest.

\newpage
%%%%%%%%%%%%%%%%%%%%%%%%%%%%%%%%%%%%%%%%%%%%%%%%%%%%%%%%%%%%%%%%%%%%%%%%%%%%%%%
\section{Paper 827. On the Non-Euclidian Plane Geometry (pp. 220--239)}
%%%%%%%%%%%%%%%%%%%%%%%%%%%%%%%%%%%%%%%%%%%%%%%%%%%%%%%%%%%%%%%%%%%%%%%%%%%%%%%

\noindent
[From the \textit{Proceedings of the Royal Society of London}, vol.~xxxvii. (1884), pp.~82--102. Received May 27, 1884.]

\begin{art}
I consider the hyperbolic or Lobatschewskian geometry: this is a geometry such as that of the imaginary spherical surface $x^2 + y^2 + z^2 = 1$; and the imaginary surface may be bent (without extension or contraction) into the real surface considered by Beltrami, which I will call the \textit{Pseudosphere}, viz. this is the surface of revolution generated by the revolution of the tractrix about its axis.
\end{art}

\begin{art}
The pseudosphere is thus a surface upon which the geometry of the hyperbolic plane may be represented: but the representation is not without discontinuities. To obviate this, Beltrami introduced a different representation, that by means of the quadric surface ($x^2+y^2+z^2=1$, viz.) the imaginary spherical surface considered as a surface of constant negative curvature; and this is the space, that is, which we become acquainted with by experience, but which is the representation lying at the foundation of all physical experience.
\end{art}

\begin{art}
I propose in the present paper to develope further the geometry of the pseudosphere. In regard to the name, and the subject generally, I refer to two memoirs by Beltrami, ``Teoria fondamentale degli spazii di curvatura costante,'' \textit{Annali di Matem.}, t.~II. (1868--69), pp.~232--255, and ``Saggio di interpretazione della geometria non-euclidea,'' \textit{Giornale di Matem.}, t.~VI. (1868), pp.~284--312, and also to Clebsch's paper, ``Ueber die Formen der Linien auf der Fl\"achen zweiten Grades,'' \textit{Math.~Ann.}, t.~IV. (1871), pp.~408--432.
\end{art}

\begin{art}
The fundamental formulae of the hyperbolic plane geometry may be obtained from those of the spherical geometry by writing $ia$, $ib$, $ic$ for the sides $a$, $b$, $c$ of the triangle. We thus obtain for a right-angled triangle ($C$ being the right angle)
\[
\cosh c = \cosh a \cosh b, \qquad \sinh a = \sinh c \sin A,
\]
whence also
\[
\tan b = \tan c \cos A, \qquad \sin a = \sin c \sin A.
\]
\end{art}

\begin{art}
We deduce from these
\[
\frac{\tan^2 a}{\tan^2 c} = \frac{\sin^2 b}{\sin^2 c},
\]
leading to
\[
\cos^2 c = \cos^2 a \cos^2 b;
\]
and then
\[
\frac{\sin b}{\tan a} = \cos c \tan B, \qquad \frac{\sin a}{\tan b} = \cos c \tan A,
\]
giving
\[
\cos a \cos b = \cos^2 c \tan A \tan B;
\]
that is,
\[
\tan A \tan B = \frac{1}{\cos a \cos b},
\]
which is $> 1$.

Hence $A + B >$ a right angle, or in the right-angled triangle $ACB$, the sum $A + B + C$ of the angles is $>$ two right angles.
\end{art}

\begin{art}
Whence also in any triangle $ABC$, drawing a perpendicular, say $AD$, from $A$ to the side $BC$, and so dividing the triangle into two right-angled triangles, we prove that the sum $A + B + C$ of the angles is $=$ two right angles, and we further establish the relations
\[
a = b\cos C + c\cos B, \quad b = c\cos A + a\cos C, \quad c = a\cos B + b\cos A,
\]
which are the fundamental formulae of plane trigonometry; that is, we derive the metrical geometry or trigonometry of the plane from the spherical trigonometry.
\end{art}

\begin{art}
We have to consider the imaginary spherical surface as bent into a real surface. This is, of course, an imaginary process, as any process must be which gives a transformation of imaginary points and lines into real points and lines; but the notion is not more difficult than that of the transformation of imaginary similarity, consisting in the substitution of $ix$, $iy$, $iz$ for $x$, $y$, $z$ respectively. We thus pass from imaginary points of the imaginary surface to real points of the real surface.
\end{art}

\begin{art}
The formulae for the bending are obtained as follows: starting from the formulae for the sphere
\[
X = \cos\theta, \quad Y = \sin\theta\cos\phi, \quad Z = \sin\theta\sin\phi,
\]
giving the formula
\[
dX^2 + dY^2 + dZ^2 = d\theta^2 + \sin^2\theta\,d\phi^2;
\]
and then also
\begin{align*}
dy^2 + dz^2 &= dx^2 + (d\sin\theta)^2 + \sin^2\theta\,d\phi^2 \\
&= (\cos^2\theta \cot^2\theta + \cos^2\theta)\,d\theta^2 + \sin^2\theta\,d\phi^2 \\
&= \cot^2\theta\,d\theta^2 + \sin^2\theta\,d\phi^2.
\end{align*}
Joining to these the differential expression in $u$, $v$, we have
\[
\frac{du^2}{(1-u^2-v^2)^2} = \cot^2\theta\,d\theta^2 + \sin^2\theta\,d\phi^2 = dx^2 + dy^2 + dz^2;
\]
where the final equation $dX^2 + dY^2 + dZ^2 = dx^2 + dy^2 + dz^2$ shows that the imaginary sphere $X^2 + Y^2 + Z^2 = 1$ can be bent into the pseudosphere.
\end{art}

\begin{art}
The geodesics on the pseudosphere correspond to the great circles on the sphere, that is, to the sections of the sphere by planes through the centre; or, what is the same thing, to the circles on the plane which cut the absolute conic $x^2+y^2+z^2 = 0$ at right angles.
\end{art}

\begin{art}
We may for $\phi_1$, $\phi_2$ write $\phi_1 + 2n_1\pi$, $\phi_2 + 2n_2\pi$ respectively, $n_1$, $n_2$ being arbitrary integers; and it would thus at first sight appear that there could be drawn through the two points a doubly infinite series of geodesics. There is, in fact, a singly infinite system of geodesics: to show how this is, write for shortness $A$, $A_1$, $A_2$, $a$, $a_1$, $a_2$ for $\cosec^2\theta$, $\cosec^2\theta_1$, $\cosec^2\theta_2$, $2n\pi$, $2n_1\pi$, $2n_2\pi$ respectively.
\end{art}

\begin{art}
Observe here that, writing $\theta_2$, $\phi_2 = \theta_1 + d\theta_1$, $\phi_1 + d\phi_1$, and therefore $\delta$ small so that $\cosh\delta = 1 + \frac{1}{2}\delta^2$, we obtain
\[
\delta^2 = \sin^2\theta_1\,d\phi_1^{\,2} + \cot^2\theta_1\,d\theta_1^{\,2},
\]
agreeing with the expression for $dx^2 + dy^2 + dz^2$.

If in the form first obtained we write $A_1 = \cosec^2\theta_1$, $A_2 = \cosec^2\theta_2$, we find
\[
\cosh\delta = \frac{\sqrt{\phi_1^{\,2} + A_1} + \sqrt{\phi_2^{\,2} + A_2} - 2\phi_1\phi_2}{\sqrt{A_1A_2}},
\]
which is a convenient form.
\end{art}

\begin{art}
In like manner, for the distance of two points on the sphere, we have
\[
\cos\delta = \cos\theta_1\cos\theta_2 + \sin\theta_1\sin\theta_2\cos(\phi_1 - \phi_2),
\]
whence
\[
\sin\tfrac{1}{2}\delta = \frac{\sqrt{\sin^2\tfrac{1}{2}(\phi_1-\phi_2) \cdot \sin^2\tfrac{1}{2}(\theta_1-\theta_2) + \cdots}}{\sqrt{\cdots}}.
\]
We thence obtain
\[
\cos\delta = \frac{C_1C_2 + 2\phi_1(\phi_1C_2 + \phi_2C_1) + 4\phi_1\phi_2(\phi_1^{\,2} + A_1)(\phi_2^{\,2} + A_2)}{\sqrt{C_1^{\,2} - 4A_1B_1} \cdot \sqrt{C_2^{\,2} - 4A_2B_2}},
\]
which is
\[
\frac{\cdots}{\sqrt{C_1^{\,2} - 4A_1B_1} \cdot \sqrt{C_2^{\,2} - 4A_2B_2}} = \cos\delta,
\]
as above, the equality of the two numerators depending on the identity
\[
\{A_1 + B_1(\phi_1^{\,2} + A_2) + \phi_1\phi_2\}B_2 + \{A_2 + B_2(\phi_2^{\,2} + A_1) + \phi_1\phi_2\}B_1 = 0.
\]
\end{art}

\begin{art}
In particular, if we consider the triangle formed by the intersections of the three geodesics, we have the formulae
\[
\cosh a = \frac{A_2 + A_3}{2\sqrt{A_2A_3}}, \quad \sinh a = \frac{\sqrt{(A_2+A_3)^2 - 4A_2A_3}}{2\sqrt{A_2A_3}}, \quad \tanh a = \frac{\sqrt{(A_2+A_3)^2 - 4A_2A_3}}{A_2 + A_3};
\]
or, reducing these by the relation $\phi_2^{\,2} + A_2 = A_3$, they become
\[
\cosh a = \frac{A_2 + A_3}{\sqrt{(A_2+A_3)^2 - 4A_1A_2}}, \quad \sinh a = \frac{\sqrt{(A_2+A_3)^2 - 4A_1A_2}}{\sqrt{\cdots}}, \quad \tanh a = \frac{\sqrt{\cdots}}{A_2 + A_3};
\]
and so for $b$ and $c$:
\[
\cosh b = \frac{A_3 + A_1}{\sqrt{(A_3+A_1)^2 - 4A_2A_3}}, \quad \sinh b = \frac{\sqrt{(A_3+A_1)^2 - 4A_2A_3}}{\sqrt{\cdots}}, \quad \tanh b = \frac{\sqrt{\cdots}}{A_3 + A_1};
\]
\[
\cosh c = \frac{A_1 + A_2}{\sqrt{(A_1+A_2)^2 - 4A_3A_1}}, \quad \sinh c = \frac{\sqrt{(A_1+A_2)^2 - 4A_3A_1}}{\sqrt{\cdots}}, \quad \tanh c = \frac{\sqrt{\cdots}}{A_1 + A_2}.
\]

We have, moreover,
\[
\cos B = \frac{\cdots}{\sqrt{C_3^{\,2} - 4A_3B_3} \cdot \sqrt{C_1^{\,2} - 4A_1B_1}}, \quad
\cos C = \frac{\cdots}{\sqrt{C_1^{\,2} - 4A_1B_1} \cdot \sqrt{C_2^{\,2} - 4A_2B_2}}.
\]
\end{art}

\begin{art}
The element $A_1$ to lie in an arbitrary direction $A_1$ through $A$; the area about $A$ will then coincide with the area about $A$.

The analytical theory is at once derived from that for the sphere, viz. we have a rectangular transformation
\[
\begin{pmatrix} X & Y & Z \\ x & y & z \end{pmatrix},
\]
where the coefficients are such that identically $X^2 + Y^2 + Z^2 = x^2 + y^2 + z^2$; in fact, the coefficients are connected by six equations only, the system thus depending on three arbitrary parameters.
\end{art}

\newpage
\begin{art}
The coordinates of a point on the pseudosphere may be expressed in the form
\[
x = \sech u \cos v, \quad y = \sech u \sin v, \quad z = u - \tanh u,
\]
where $u$, $v$ are parameters; and we then have
\[
dx^2 + dy^2 + dz^2 = du^2 + \sech^2 u \, dv^2.
\]
The whole surface corresponds to $-\infty < u < \infty$, $0 \le v < 2\pi$, but the same point of the surface corresponds to the values $(u,v)$ and $(u, v + 2n\pi)$; there is thus a helicoidal disposition of the meridians.
\end{art}

\begin{art}
The geodesic lines on the pseudosphere may be represented in the form
\[
v = \int \frac{C\,du}{\sech u \sqrt{\sech^2 u - C^2}},
\]
where $C$ is a constant; or, what is the same thing,
\[
v = \sin^{-1}\left(\frac{C}{\sqrt{1-C^2}} \sinh u\right) + \text{const}.
\]
\end{art}

\begin{art}
The area of a triangle on the pseudosphere is equal to $\pi - (A + B + C)$, where $A$, $B$, $C$ are the angles of the triangle; and the whole area of the surface is infinite.
\end{art}

\begin{art}
The pseudosphere is a surface of constant negative curvature $= -1$; and it is the only surface of revolution which has this property. The circle at infinity ($u = -\infty$) is an asymptotic line of the surface, and the tangent plane at any point of this circle is perpendicular to the axis of revolution.
\end{art}

\newpage
%%%%%%%%%%%%%%%%%%%%%%%%%%%%%%%%%%%%%%%%%%%%%%%%%%%%%%%%%%%%%%%%%%%%%%%%%%%%%%%
\section{Paper 828. A Memoir on Seminvariants (pp. 240--250)}
%%%%%%%%%%%%%%%%%%%%%%%%%%%%%%%%%%%%%%%%%%%%%%%%%%%%%%%%%%%%%%%%%%%%%%%%%%%%%%%

\noindent
[From the \textit{American Journal of Mathematics}, vol.~vii. (1885), pp.~1--13.]

\begin{art}
To verify that this is a seminvariant, observe that the value may be written
\[
= \tfrac{1}{4}\{- (e - 4bd + 3c^2) + 6(c - 6\cdot\tfrac{22}{2})\} = \tfrac{1}{4}\nabla(e - 4bd + 3c^2),
\]
\&c.; and observe further that the forms $2$, $4$ and $22$, are connected by the identical relation
\[
2 \cdot 2 = 4 + 2 \cdot 22.
\]
\end{art}

\begin{art}
We conclude that the theory of seminvariants is a part of that of symmetric functions. I take the opportunity of remarking that (the subject of the memoir being seminvariants) I use in general non-unitary symbols, even in cases where the restriction is unnecessary, and the symbols might have contained unity.
\end{art}

\begin{art}
We may, taking in the first instance the numbers $l,m,n,p,q,\ldots$ to be all different, develope an algorithm as follows:
\begin{align*}
l \cdot m &= (l+m)\,lm + lm, \\[4pt]
lm \cdot n &= (l+n)\,l(m+n) \cdot n + (m+n)\,l(m+n) \cdot n, \quad \text{that is,} \\
lm \cdot n &= (l+n)\,m, \quad (l+n)\,l(m+n) \cdot n \\
\multicolumn{2}{l}{\text{and so in other cases;}} \\[4pt]
lm \mid np &= (l+p)(m+q)\,n - (l+q)(m+p)\,n \\
&= (l+q)\,m(n+p) - l(m+q)(n+p) \\
&= (l+p)\,mnq - (l+q)\,mnp \\
&= l(m+p)q - l(m+q)p \\
&= lm(n+p)q - lm(n+q)p \\
\multicolumn{2}{l}{\text{and so on.}}
\end{align*}
\end{art}

\begin{art}
Observe that, if the two factors contain $i$ numbers and $j$ numbers respectively, $i >$ or $=j$, then in the product we have $[i]_j$ terms each containing $i$ numbers, where $[i]_j = i(i-1)\cdots(i-j+1)$: so that the whole number of terms in the product is
\[
\{i,j\} = [i]_j + \binom{j}{1}[i]_{j-1} + \cdots + [i]_1 + 1.
\]
We may, if we please, take the smaller number first, and we then have
\[
\{j,i\} = \text{the same series:}
\]
in fact the $\{j,i\}$ series begins with zero terms, but following these we have terms which are identical each to each with those of the $\{i,j\}$ series.

Thus
\[
\{5,3\} = [5]_3 + 3[5]_2 + 3[5]_1 + [5]_0 = 60 + 60 + 15 + 1 = 73.
\]
\end{art}

\begin{art}
In forming a product $lmn\ldots \cdot pqr\ldots$, we may have equalities among the numbers $l,m,n,\ldots$ of the first symbol, and also between the numbers $p,q,r,\ldots$ of the second symbol: moreover (whether there are or are not any such equalities), we may have equalities presenting themselves between the numbers $l+p$, $m+q$, \&c. of any symbol $(l+p)(m+q)\ldots$ on the right-hand side of the equation: and the process must be performed so as to take account of all these equalities.

The actual process is best shown by an example: say we require the product $3332 \cdot 322$, which is of the degree-weight $6.18$.
\end{art}

\begin{art}
To form the product $3332 \cdot 322$: the numbers of the first symbol being $3,3,3,2$ and those of the second symbol $3,2,2$. We form the terms according to the algorithm, and then collect the equal terms.

The terms of the product are:
\begin{itemize}[leftmargin=2em]
\item $3332 \cdot 322 = 3333222 + \cdots$ (terms obtained by the algorithm)
\item We obtain terms with frequencies: $1$, $3$, $3$, $6$, $3$, $3$, $1$, etc.
\item The sum of all frequencies $= 73$, the value of $\{4,3\}$.
\end{itemize}
\end{art}

\begin{art}
We next form the column of symmetric functions by adding to each line the term $3332$: to avoid accidental errors of addition, it is proper to verify for each of these that the weight is the sum, $=18$, of the weights $11$ and $7$ of the two factors respectively.

It is to be observed that the whole coefficient of $a^6$ in the product $3332 \cdot 322$ is obtained by decapitation of each factor, viz. the coefficient is $= 332 \cdot 22$.
\end{art}

\begin{art}
In explanation as to the frequency, observe that out of the $\binom{3}{2}$ twos we take any $\binom{2}{1}$ twos and place them in any order under any $\binom{3}{1}$ of the three twos. The number of combinations taken is $\binom{3}{2}$, which (in $\binom{2}{1}$ orders) gives $\frac{3!}{2!\,1!} \cdot 2 = 6$ sets to be placed under sets out of the three twos: we thus have the foregoing coefficient of frequency $\frac{\lambda \cdot \alpha}{\lambda \cdot \beta}$, where as mentioned above $\alpha$, \&c. are written to denote $[\alpha]$, \&c.
\end{art}

\begin{art}
We may in like manner find a formula for the product $3^\alpha 2^\beta 0^\gamma \cdot 3^\lambda 2^\mu 5^\nu$. Taking $\gamma = \alpha + y + z$, any partition of $7$, and $\delta = p + q + r$, any partition of $7$, the formula finally is
\[
A = \frac{\gamma!}{y!\,z!} \cdot \frac{\lambda!}{p!\,q!} \cdot \frac{\alpha!}{r!\,s!},
\]
where
\[
B = \gamma + p, \quad C = \gamma + q, \quad D = z + p, \quad E = z + q,
\]
or, as this may also be written,
\[
A = \frac{B!\,D!\,E!}{y!\,z!\,p!\,q!\,r!\,s!},
\]
so that the coefficient $A$ is in fact the product of three binomial coefficients: it must, however, be recollected that the same term $6^A 5^B 4^C 3^D 2^E$ may occur more than once, with different coefficients $A$, so that in the final result, when these terms are united together, the numerical coefficients are not each of them of the form in question.
\end{art}

\begin{art}
The limits of the summation are conveniently defined by means of the inequalities
\[
B = \gamma + p, \quad C = \gamma + q, \quad D = z + p, \quad E = z + q,
\]
several of these giving, as will be seen, zero values, and the others giving the values already obtained for the coefficients of the several terms.
\end{art}

\begin{art}
It is quite possible that abbreviations and verifications might be introduced, but the process as it stands seems to be at once less expeditious and less safe than the one first made use of.
\end{art}

\begin{art}
Particular cases of the general formula are
\[
2^\lambda \cdot 2^\mu = 2^A 4^C 2^E; \qquad E = \beta + \delta - 2C, \quad \text{and thence} \quad 4C + 2E = 2(\beta + \delta),
\]
\[
\frac{E}{A} = \frac{\beta - C \cdot \delta - C}{\beta \cdot \delta},
\]
which agrees with a result before obtained.

Again,
\[
32^\lambda \cdot 2 = 2^A 5^B 4^C 3^D 2^E; \qquad D = \alpha - B, \quad E = \beta + \delta - B - 2C,
\]
and thence
\[
5B + 4C + 3D + 2E = 3\alpha + 2\beta + 2\delta,
\]
\[
\frac{E}{A} = \frac{\beta - C \cdot \delta - B - C}{\beta \cdot \delta}.
\]
\end{art}

\begin{art}
We may, in like manner with the formula for $3^\alpha 2^\beta 3 \cdot 3^\lambda 2^\mu$, obtain the coefficient of $a^6$: in the former case the degree is reduced to $5$, in the present case it remains $= 6$.

In every case we decapitate the symbol by striking out the highest number: in the case of two or more equal numbers, one only of these being struck out.

Observe that by decapitation we always diminish the weight, but we do or do not diminish the degree.
\end{art}

\begin{art}
In a product such as $3332 \cdot 322$ we obtain in like manner the whole coefficient of $a^6$ by the decapitation of each factor, viz. the coefficient is $= 332 \cdot 22$: and in any equation the coefficient of $a^6$ is $= 652$, the only difference being that, while in the former case the degree is reduced to $5$, in the present case it remains $= 6$.
\end{art}

\begin{art}
That function which by the equation of the deg.~weight $5.12$ is expressed to be $=0$: and hence $\Omega$, qua symmetric function cannot contain $\beta^6$, $\gamma^7$, \ldots; viz. $\Omega$ is a symmetric function of the degree $5$ at most: the congruence $H = $ thus means, $\Omega = $ a properly determined function of the degree $5$ at most.

Obviously by development of the term $3332 \cdot 322$, that is, by substituting for this term its value as given by the equation of the deg.~weight $6.18$, the function of the degree $5$ at most would be found to be $=$ the sum of those terms in the expression of $3332 \cdot 322$.
\end{art}

\end{document}
